\section{问题二：基于能力互补性的场景效用评价}

\subsection{问题分析}

问题一得到的是模型在统一公开评测体系中的基础能力画像，而问题二关注这些能力进入具体任务环境后能够产生的使用效用。二者的关系可概括为
\[
\text{基础能力评价}\longrightarrow\text{场景效用评价}.
\]
因此，本文不在问题二中重新构建 Benchmark 指标体系，而是继承问题一冻结后的五维能力向量，并分别面向科研长文本分析、大众日常通用对话和计算机代码开发建立差异化评价模型。

普通线性加权模型允许不同能力完全补偿，即一项能力的极高得分可以抵消另一项关键能力的严重不足。然而科研长文本分析需要长上下文、复杂推理与事实可靠性协同，代码开发也依赖代码工程能力与复杂推理共同发挥作用。为同时刻画场景对能力重要性和可替代性的影响，本文构造“最小信息偏移场景权重+CES互补型效用”模型，简称 KL--CES 模型。其核心思想是：场景不仅改变各能力的重要程度，还改变能力之间的可替代结构。

\subsection{基础能力与场景需求映射}

设模型 $i$ 在能力 $j$ 上由问题一给出的最终得分为 $z_{ij}$。五个能力维度依次为复杂推理能力 $C_1$、知识与事实可靠性 $C_2$、长上下文处理能力 $C_3$、代码与软件工程能力 $C_4$ 和多模态能力 $C_5$。若问题一输出采用 $[0,100]$ 量纲，则先作归一化
\[
\widetilde z_{ij}=\frac{z_{ij}}{100},\qquad \widetilde z_{ij}\in[0,1].
\]

科研长文本分析的核心能力链为 $C_3\rightarrow C_1\rightarrow C_2$；大众日常通用对话主要依赖 $C_2,C_1,C_5$，并由 $C_3$ 支持上下文连续性；计算机代码开发以 $C_4$ 和 $C_1$ 为核心，同时受到事实可靠性和上下文处理能力的辅助影响。价格、时延、算力与能耗等成本变量留在问题三的性能--成本决策层中处理，以避免评价层次重叠。

现有冻结数据中的 Agent/Tool 指标未形成覆盖全部模型且协议一致的独立矩阵，并且相关软件工程 Agent 指标已部分进入 $C_4$。因此本文不额外设置工具调用维度，而由 $C_4$ 与 $C_1$ 共同表征代码开发中的主要 Agent 型需求。

\subsection{基于最小信息偏移的场景权重修正}

设问题一提供的客观信息权重为
\[
\boldsymbol w^{(0)}=(w_1^{(0)},\ldots,w_m^{(0)}),
\qquad w_j^{(0)}>0,\quad \sum_{j=1}^{m}w_j^{(0)}=1.
\]
该权重反映五类能力在现有公开评测数据中的信息贡献、非冗余性和稳定性，是缺乏具体场景偏好时的中性参考分布，而不是现实场景中的直接价值权重。

对场景 $s$，设置权重 $\boldsymbol w^{(s)}$，并通过需求约束集合 $\Omega_s$ 描述核心能力总权重、能力下界及重要性顺序。场景权重由下式确定：
\[
\boxed{
\boldsymbol w^{(s)}=
\arg\min_{\boldsymbol w\in\Omega_s}
D_{\mathrm{KL}}\!\left(\boldsymbol w\Vert\boldsymbol w^{(0)}\right)
}
\]
其中
\[
D_{\mathrm{KL}}\!\left(\boldsymbol w\Vert\boldsymbol w^{(0)}\right)
=\sum_{j=1}^{m}w_j\ln\frac{w_j}{w_j^{(0)}}.
\]
公共约束为 $w_j\ge 0$ 和 $\sum_jw_j=1$。三类场景的核心约束写为
\begin{align*}
w_{C_1}^{(R)}+w_{C_2}^{(R)}+w_{C_3}^{(R)}&\ge\alpha_R,\\
w_{C_1}^{(D)}+w_{C_2}^{(D)}+w_{C_5}^{(D)}&\ge\alpha_D,\\
w_{C_4}^{(C)}+w_{C_1}^{(C)}&\ge\alpha_C.
\end{align*}
此外设置若干参数化顺序或下界约束，例如科研场景中 $w_{C_3}\ge w_{C_4}$、$w_{C_1}\ge w_{C_5}$，代码场景中对 $w_{C_4}$ 设置最低需求。所有阈值均作为敏感性参数而非预先认定的唯一真值。

为检验场景排名是否过度依赖问题一客观权重先验，本文另取等权参考
\[
\boldsymbol w^{(\mathrm{eq})}=(1/m,\ldots,1/m),
\]
在相同约束下重复 KL 投影，比较两类先验下的排名稳健性。

\subsection{基于CES函数的能力互补型场景效用模型}

在场景权重确定后，采用常替代弹性函数刻画能力间的不完全替代关系：
\[
\boxed{
U_i^{(s)}=
\left[
\sum_{j=1}^{m}w_j^{(s)}
(\widetilde z_{ij}+\varepsilon)^{\rho_s}
\right]^{1/\rho_s}
}
\]
其中 $\varepsilon>0$ 为统一设置的数值稳定常数，$\rho_s$ 控制能力之间的替代程度，相应替代弹性为
\[
\sigma_s=\frac{1}{1-\rho_s}.
\]
当 $\rho_s\rightarrow1$ 时，CES 接近线性加权；当 $\rho_s\rightarrow0$ 时，CES 连续收敛到加权几何平均
\[
U_i^{(s)}=\prod_{j=1}^{m}(\widetilde z_{ij}+\varepsilon)^{w_j^{(s)}};
\]
当 $\rho_s<0$ 时，关键能力短板受到更强惩罚。科研与代码场景通常具有更明显的互补性，而日常对话的能力替代性相对较高。本文不把单个 $\rho_s$ 当作已知真值，而是在预设合理区间内进行参数敏感性分析。

为检验互补结构的增量作用，同时构造完全补偿型线性基准
\[
L_i^{(s)}=\sum_{j=1}^{m}w_j^{(s)}\widetilde z_{ij}.
\]
线性模型仅作为基准，KL--CES 仍为问题二的主模型。通过比较 $\operatorname{Rank}_{\mathrm{CES}}$ 与 $\operatorname{Rank}_{\mathrm{Linear}}$，可以识别能力失衡对场景评价产生的影响。

\subsection{能力失衡惩罚指数}

定义模型 $i$ 在场景 $s$ 中因能力不均衡产生的绝对效用损失为
\[
P_i^{(s)}=L_i^{(s)}-U_i^{(s)},
\]
并定义能力失衡惩罚指数（Performance Imbalance Penalty Index, PRI）
\[
\boxed{
\mathrm{PRI}_i^{(s)}=
\frac{L_i^{(s)}-U_i^{(s)}}{L_i^{(s)}}
}.
\]
PRI 衡量模型在互补型场景下相对于完全补偿型评价所损失的效用比例。较大的 PRI 表明模型存在较明显的关键能力短板；其解释必须与所采用的 $\rho_s$ 区间共同报告。

\subsection{场景排名迁移与作用机理分析}

设问题一最终综合排名为 $R_i^{(0)}$，场景 $s$ 的 KL--CES 排名为 $R_i^{(s)}$，定义
\[
\boxed{\Delta R_i^{(s)}=R_i^{(s)}-R_i^{(0)}}.
\]
由于名次 1 最优，$\Delta R_i^{(s)}<0$ 表示进入该场景后排名上升，$\Delta R_i^{(s)}>0$ 表示排名下降。若问题一最终综合排名不适合作为统一基线，则只绘制三个场景之间的 Bump Chart，并保持场景模型本身不变。

排名迁移的机理由两条路径共同解释：其一是 KL 投影改变能力的重要性结构 $\boldsymbol w^{(s)}$；其二是 CES 参数改变能力的可替代结构 $\rho_s$。成对竞争分析进一步比较相邻模型的原始能力差、线性加权能力效应、CES 曲率效应、能力失衡惩罚以及优先改善方向。

\subsection{模型边际能力贡献分析}

由于 CES 为非线性函数，不能直接把 $w_jz_{ij}$ 称为 CES 能力贡献。本文采用两种互补解释量。第一种为 CES 内层聚合份额
\[
q_{ij}^{(s)}=
\frac{w_j^{(s)}(\widetilde z_{ij}+\varepsilon)^{\rho_s}}
{\sum_k w_k^{(s)}(\widetilde z_{ik}+\varepsilon)^{\rho_s}},
\qquad \sum_jq_{ij}^{(s)}=1.
\]
它描述各能力在 CES 内部聚合项中的相对份额，但不等同于可加总的最终效用贡献。

第二种为边际效用。记
\[
A_i^{(s)}=\sum_jw_j^{(s)}(\widetilde z_{ij}+\varepsilon)^{\rho_s},
\]
则
\[
\frac{\partial U_i^{(s)}}{\partial \widetilde z_{ij}}
=w_j^{(s)}(\widetilde z_{ij}+\varepsilon)^{\rho_s-1}
\left(U_i^{(s)}\right)^{1-\rho_s}.
\]
进一步定义效用弹性
\[
E_{ij}^{(s)}=
\frac{\partial U_i^{(s)}}{\partial \widetilde z_{ij}}
\frac{\widetilde z_{ij}}{U_i^{(s)}}.
\]
边际效用用于回答“改进该模型的哪项能力最能提高当前场景效用”。程序中同时采用数值微分检验解析导数的一致性。

\subsection{模型敏感性与稳健性分析}

场景约束与替代参数均具有一定主观性，因此在每个场景内构造二维参数网格
\[
(\alpha_s,\rho_s)\in
[\alpha_s^L,\alpha_s^U]\times[\rho_s^L,\rho_s^U],
\]
对每个网格点重新求解 KL 权重并计算 CES 排名。对每个模型统计平均排名、最好排名、最差排名、Top3 比例、排名标准差和相邻参数点的排名突变，以识别稳定区域与阈值效应。同时在问题一客观先验和等权先验下重复计算，以检验先验依赖性。

针对能力缺失，本文不采用填 0。根据问题一最终交付情况，分别使用完整案例、共同维度、问题一冻结策略继承或区间传播。共同维度方案会在所有模型共同可观测的维度集合上重新求解场景约束权重，并明确标记结果口径；区间传播则计算效用上下界以及最好和最差可能排名。

% TODO: Insert final Q1-dependent results after Q1 freeze
